\scanpage{8}

% archive sheet number: 7

\margin{\unsure{2|$Z$|$\big\}$}}

Mais attention, ce résultat ne doit pas faire illusion : Si on se place dans la catégorie des \uline{Schémas} sur $\operatorname{Spec}(A)$, il n'implique pas l'existence du quotient dans cette catégorie, \deleted{\unsure{ni}} et si le quotient y existe, il peut être distinct de celui qu'on vient de définir.

Supposons p.ex $A={}$corps alg. clos $k$, $\deleted{\unsure{E}}C={}$algèbre \inserted{affine} d'un groupe linéaire connexe non \unsure{résoluble} $\deleted{\unsure{G}}E$, $B={}$algèbre d'un sous-groupe de Borel $\deleted{\unsure{B}}G$. Alors $E/G$ dans la catégorie des schémas existe bien et est une variété complète $\ne e$. Il en résulte que le quotient dans la \uline{catégorie des variétés affines} est réduit au point $e$.

On \uline{doit donc vérifier} \uline{qu'en l'occurrence}, \deleted{\unsure{les}} \makebox[0pt][r]{\smash{\raisebox{-.9\baselineskip}{\(\left\Vert\vphantom{\rule{0pt}{2.2\baselineskip}}\right.\)}}\hspace{.6em}}[\uline{$C$ étant principal homogène}] \deleted{les éléments} \uline{$C'$ est isomorphe \unsure{:|à} $A$} \uline{[par $\lambda\to\lambda1_C$]}.

En effet, si $x$ est semi-primitif, alors le diagramme commutatif $(*)$ donne

\[
x\otimes1_C=1_C\otimes x\qquad\text{dans}\quad C\otimes C
\]
